% ══════════════════════════════════════════════════════════════════════
% Volume IV-B, Chapters 1--5
% ══════════════════════════════════════════════════════════════════════

% ══════════════════════════════════════════════════════════════════════
% Chapter 1: The Decay of Meaning
% ══════════════════════════════════════════════════════════════════════

\chapter{The Decay of Meaning: Why Reinterpretation Always Loses}
\label{ch:decay}

\epigraph{``A word is dead / When it is said, / Some say. /
I say it just / Begins to live / That day.''}{Emily Dickinson}

% ── Historical Preamble ──────────────────────────────────────────────

The childhood game of ``telephone''---in which a message is whispered
from ear to ear around a circle, arriving at the end comically
garbled---is not merely a game.  It is a \emph{theorem}.

In this chapter, we prove that every act of reinterpretation---every
composition of a sign with a new context---introduces an irreversible
distortion.  The distortion is not random noise; it has a precise
algebraic structure.  It decomposes into a \emph{bias term} (the
context's own resonance with the observer) and an \emph{emergence term}
(the nonlinear interaction between sign and context).  Only when both
are zero---only when the context is silence---is the reinterpretation
lossless.

This result overturns the naive hope that ``careful'' reinterpretation
can preserve meaning perfectly.  It cannot.  Every reading leaves a mark
on the text, and the marks accumulate.

\section{Sign Chains and Reinterpretation}

We begin by formalizing the notion of a \emph{sign chain}: a sequence
of reinterpretations, each modeled as composition with a new context.

\begin{definition}[Sign Chain]
\label{def:sign-chain}
Let $a \in \IS$ be a sign and $c_1, c_2, \ldots, c_n \in \IS$ be
a sequence of contexts.  The \emph{$n$-step sign chain} is defined
recursively:
\[
  \mathrm{chain}(a, []) = a, \qquad
  \mathrm{chain}(a, c :: \mathrm{rest}) = \mathrm{chain}(a \compose c, \mathrm{rest}).
\]
\end{definition}

\begin{lstlisting}
def signChain (a : I) : List I → I
  | [] => a
  | c :: rest => signChain (compose a c) rest
\end{lstlisting}

The definition is straightforward, but its consequences are profound.

\begin{theorem}[Sign Weight Growth]
\label{thm:sign-weight-growth}
For any sign $a$ and context $c$,
\[
  \rs(a \compose c, a \compose c) \geq \rs(a, a).
\]
Every reinterpretation increases or preserves the sign's weight.
\end{theorem}

\begin{proof}
This is an immediate consequence of axiom E4 (compositional enrichment):
$\rs(a \compose b, a \compose b) \geq \rs(a, a)$ for all $a, b$.
\end{proof}

\begin{remark}[Why this is surprising]
At first glance, this seems like good news: reinterpretation makes signs
``heavier.''  But weight measures \emph{quantity of presence}, not
\emph{fidelity to the original meaning}.  A sign can grow heavier while
drifting arbitrarily far from its original resonance profile.  The sign
gets more ``substantial'' while becoming less ``faithful.''  This is the
paradox of tradition: cultural transmission accumulates weight while
potentially distorting meaning.
\end{remark}

\section{The Reinterpretation Gap}

\begin{definition}[Reinterpretation Gap]
The \emph{reinterpretation gap} measures how much one step of
reinterpretation changes the resonance of a sign with an observer:
\[
  \mathrm{gap}(a, c, \mathrm{obs}) := \rs(a \compose c, \mathrm{obs}) - \rs(a, \mathrm{obs}).
\]
\end{definition}

\begin{lstlisting}
noncomputable def reinterpretationGap (a c observer : I) : R :=
  rs (compose a c) observer - rs a observer
\end{lstlisting}

\begin{theorem}[Gap Decomposition]
\label{thm:gap-decomp}
The reinterpretation gap decomposes as:
\[
  \mathrm{gap}(a, c, \mathrm{obs}) = \rs(c, \mathrm{obs}) + \emerge(a, c, \mathrm{obs}).
\]
\end{theorem}

\begin{proof}
By the resonance composition equation $\rs(a \compose c, \mathrm{obs})
= \rs(a, \mathrm{obs}) + \rs(c, \mathrm{obs}) + \emerge(a, c, \mathrm{obs})$,
subtraction gives the result immediately.
\end{proof}

\begin{lstlisting}
theorem reinterpretationGap_decomp (a c observer : I) :
    reinterpretationGap a c observer =
    rs c observer + emergence a c observer := by
  unfold reinterpretationGap
  have h := rs_compose_eq a c observer
  linarith
\end{lstlisting}

\begin{interpretation}[The two channels of distortion]
The gap has two independent components:
\begin{enumerate}
\item The \textbf{bias term} $\rs(c, \mathrm{obs})$: the context's own
  resonance with the observer.  This is the ``literal'' contribution of
  the context---what it brings to the table independently of the sign.
\item The \textbf{emergence term} $\emerge(a, c, \mathrm{obs})$: the
  nonlinear interaction between sign and context.  This is the ``creative''
  distortion---the new meaning that arises from the composition.
\end{enumerate}
Both terms are generally nonzero.  To achieve zero gap, \emph{both}
must vanish simultaneously.  This is a much stronger condition than
either alone.
\end{interpretation}

\section{Two-Step Distortion: The Telephone Theorem}

\begin{theorem}[Two-Step Distortion]
\label{thm:two-step}
When a sign is reinterpreted through two successive contexts
$c_1, c_2$, the total gap equals the sum of individual gaps
\emph{plus} a cross-emergence term:
\[
  \rs(a \compose c_1 \compose c_2, \mathrm{obs}) - \rs(a, \mathrm{obs})
  = \mathrm{gap}(a, c_1, \mathrm{obs}) +
    \mathrm{gap}(a \compose c_1, c_2, \mathrm{obs}).
\]
\end{theorem}

\begin{proof}
This is a telescoping identity:
\begin{align*}
  &\rs((a \compose c_1) \compose c_2, \mathrm{obs}) - \rs(a, \mathrm{obs}) \\
  &= [\rs((a \compose c_1) \compose c_2, \mathrm{obs})
    - \rs(a \compose c_1, \mathrm{obs})] +
    [\rs(a \compose c_1, \mathrm{obs}) - \rs(a, \mathrm{obs})] \\
  &= \mathrm{gap}(a \compose c_1, c_2, \mathrm{obs}) +
    \mathrm{gap}(a, c_1, \mathrm{obs}).
\end{align*}
\end{proof}

\begin{lstlisting}
theorem two_step_distortion (a c1 c2 observer : I) :
    reinterpretationGap a c1 observer +
    reinterpretationGap (compose a c1) c2 observer =
    rs (compose (compose a c1) c2) observer - rs a observer := by
  unfold reinterpretationGap; ring
\end{lstlisting}

\begin{remark}[Why this is surprising]
The naive expectation is that two ``small'' reinterpretations should
produce a ``small'' total distortion.  But each gap decomposes into
bias + emergence, and the \emph{emergence terms are generally correlated}
via the cocycle condition.  Two individually small gaps can produce a
large total distortion if the cross-emergence aligns.  Distortion is
not additive---it is \emph{super-additive} when emergence cooperates.
\end{remark}

\section{The Only Lossless Reinterpretation}

\begin{theorem}[Void Preserves Perfectly]
\label{thm:void-preserves}
The only context that produces zero reinterpretation gap for all
observers is the void:
\[
  \mathrm{gap}(a, \void, \mathrm{obs}) = 0 \quad \text{for all } a, \mathrm{obs}.
\]
\end{theorem}

\begin{proof}
By axiom A3, $a \compose \void = a$, so
$\rs(a \compose \void, \mathrm{obs}) = \rs(a, \mathrm{obs})$
and the gap is zero.
\end{proof}

\begin{lstlisting}
theorem void_reinterpretation_zero (a observer : I) :
    reinterpretationGap a (void : I) observer = 0 := by
  unfold reinterpretationGap; simp [void_right']
\end{lstlisting}

\begin{interpretation}[Silence is the only perfect reading]
This theorem says that the \emph{only} way to ``reinterpret'' a text
without changing it is to add nothing.  Every substantive reading---every
context that is not pure silence---necessarily changes what the text
means to some observer.  This is the formal version of reader-response
theory's central claim: the reader always contributes.  And the only
reader who contributes nothing is the one who doesn't read.
\end{interpretation}

\section{Non-Void Signs Leave Permanent Marks}

\begin{theorem}[Irreversibility of Reading]
\label{thm:reading-irreversible}
If $a \neq \void$, then for any context $c$:
\[
  \rs(a \compose c, a \compose c) > 0.
\]
The reinterpreted sign always has strictly positive weight.
\end{theorem}

\begin{proof}
Since $a \neq \void$, we have $\rs(a, a) > 0$ by the nondegeneracy axiom.
By compositional enrichment, $\rs(a \compose c, a \compose c) \geq \rs(a, a) > 0$.
\end{proof}

\begin{lstlisting}
theorem nonvoid_reinterpretation_positive (a c : I)
    (ha : a ≠ void) :
    rs (compose a c) (compose a c) > 0 := by
  have h1 := compose_enriches' a c
  have h2 := rs_self_pos a ha
  linarith
\end{lstlisting}

\begin{remark}[Why this is surprising]
This means that once an idea enters the cultural conversation, it can
never be fully erased.  Composition with any context---even a hostile
one, even an attempt at refutation---produces something with positive
weight.  Censorship, in the ideatic-space sense, is impossible: engaging
with an idea to suppress it only adds to its weight.  The only way to
truly silence an idea is to never engage with it at all (compose with
void).
\end{remark}


% ══════════════════════════════════════════════════════════════════════
% Chapter 2: The Impossibility of Perfect Translation
% ══════════════════════════════════════════════════════════════════════

\chapter{The Impossibility of Perfect Translation}
\label{ch:translation}

\epigraph{``Traduttore, traditore.''\\(Translator, traitor.)}{Italian proverb}

Translation has been called ``the art of failure'' (Umberto Eco).  In this
chapter, we prove that this is not merely a literary sentiment but a
\emph{mathematical theorem}.  Any translation that is not the identity
must introduce distortion.  The distortion is measurable, it accumulates
through chains of retranslation, and it has a precise lower bound.

\section{Translation Loss: The Basic Inequality}

\begin{definition}[Translation Loss]
For a translation map $\varphi : \IS \to \IS$, the \emph{translation
loss} for a pair $(a, b)$ is:
\[
  \mathrm{loss}(\varphi, a, b) := |\rs(\varphi(a), \varphi(b)) - \rs(a, b)|.
\]
\end{definition}

\begin{lstlisting}
noncomputable def translationLoss (phi : I → I) (a b : I) : R :=
  |rs (phi a) (phi b) - rs a b|
\end{lstlisting}

\begin{theorem}[Faithful Translations Have Zero Loss]
\label{thm:faithful-zero}
If $\varphi$ is faithful (preserves all resonances), then
$\mathrm{loss}(\varphi, a, b) = 0$ for all $a, b$.
\end{theorem}

\begin{proof}
If $\rs(\varphi(a), \varphi(b)) = \rs(a, b)$ for all $a, b$,
then the absolute value of their difference is zero.
\end{proof}

\begin{theorem}[The Self-Distortion Lower Bound]
\label{thm:self-distortion}
If a translation changes the self-resonance of even \emph{one} idea,
then the translation loss is strictly positive for that idea:
\[
  \rs(\varphi(a), \varphi(a)) \neq \rs(a, a)
  \implies
  \mathrm{loss}(\varphi, a, a) > 0.
\]
\end{theorem}

\begin{proof}
If $\rs(\varphi(a), \varphi(a)) \neq \rs(a, a)$, then their difference
is nonzero, so its absolute value is positive.
\end{proof}

\begin{lstlisting}
theorem translation_self_loss (phi : I → I) (a : I)
    (h : rs (phi a) (phi a) ≠ rs a a) :
    translationLoss phi a a > 0 := by
  unfold translationLoss
  exact abs_pos.mpr (sub_ne_zero.mpr h)
\end{lstlisting}

\begin{remark}[Why this is surprising]
The theorem says that you cannot translate ``approximately faithfully.''
If a translation preserves \emph{all} resonances except one self-resonance,
it is still unfaithful, and the unfaithfulness is \emph{detectable}:
the translation loss is strictly positive.  There is no ``good enough''
approximation---faithfulness is all or nothing.
\end{remark}

\section{Chain Translation: Distortion Accumulates}

\begin{theorem}[Chain Distortion Decomposition]
\label{thm:chain-decomp}
Composing two translations $\psi \circ \varphi$ produces distortion
that decomposes into two layers:
\[
  \rs(\psi(\varphi(a)), \psi(\varphi(b))) - \rs(a, b)
  = \underbrace{[\rs(\psi(\varphi(a)), \psi(\varphi(b)))
    - \rs(\varphi(a), \varphi(b))]}_{\text{second layer}} +
    \underbrace{[\rs(\varphi(a), \varphi(b)) - \rs(a, b)]}_{\text{first layer}}.
\]
\end{theorem}

\begin{proof}
This is a telescoping identity; it follows from algebraic rearrangement.
\end{proof}

\begin{lstlisting}
theorem chain_translation_loss_decomp (phi psi : I → I)
    (a b : I) :
    rs (psi (phi a)) (psi (phi b)) - rs a b =
    (rs (psi (phi a)) (psi (phi b)) - rs (phi a) (phi b)) +
    (rs (phi a) (phi b) - rs a b) := by ring
\end{lstlisting}

\begin{interpretation}[Why retranslation degrades]
Each layer of translation adds its own distortion.  Even if each individual
translation is ``small'' (low distortion), the total can be large because
the layers are \emph{independent}.  This is the formal reason that
translating a poem from Chinese to French to English produces a worse
result than translating directly from Chinese to English: each intermediate
step adds its own layer of loss, and the losses add.
\end{interpretation}

\section{Round-Trip Loss: The Myth of Perfect Back-Translation}

\begin{theorem}[Round-Trip Loss]
\label{thm:round-trip}
For any two translations $\varphi, \psi$, the round-trip distortion
satisfies:
\[
  \rs(\psi(\varphi(a)), \psi(\varphi(a))) - \rs(a, a)
  = \underbrace{[\rs(\psi(\varphi(a)), \psi(\varphi(a)))
    - \rs(\varphi(a), \varphi(a))]}_{\text{back-translation loss}} +
    \underbrace{[\rs(\varphi(a), \varphi(a)) - \rs(a, a)]}_{\text{forward loss}}.
\]
\end{theorem}

\begin{proof}
Again a telescoping identity.
\end{proof}

\begin{lstlisting}
theorem roundtrip_loss (phi psi : I → I) (a : I) :
    rs (psi (phi a)) (psi (phi a)) - rs a a =
    (rs (psi (phi a)) (psi (phi a)) - rs (phi a) (phi a)) +
    (rs (phi a) (phi a) - rs a a) := by ring
\end{lstlisting}

\begin{remark}[Why this is surprising]
The theorem destroys the comforting notion of ``back-translation as
quality check.''  Even if you translate and back-translate, the round-trip
distortion is the \emph{sum} of forward and backward losses.  Having a
``back-translation'' does not guarantee zero net loss.  The round trip
through Chinese $\to$ French $\to$ Chinese loses meaning in \emph{both}
directions.
\end{remark}

\section{Void-Preservation Is Not Sufficient}

\begin{theorem}[Void-Preserving Is Not Faithful]
\label{thm:void-not-faithful}
A void-preserving translation ($\varphi(\void) = \void$) can still
have arbitrarily high translation loss.  Void-preservation is necessary
but not sufficient for faithfulness.
\end{theorem}

\begin{proof}
We prove a disjunction: either $\varphi$ is faithful (zero loss everywhere),
or there exist $a, b$ with positive loss.  If $\varphi$ is not faithful,
there exist $a, b$ with $\rs(\varphi(a), \varphi(b)) \neq \rs(a, b)$,
giving positive loss.
\end{proof}

\begin{lstlisting}
theorem void_preserving_not_sufficient :
    ∀ phi : I → I, ds_voidPreserving phi →
    (ds_faithful phi ∨ ∃ a b, translationLoss phi a b > 0) := by
  intro phi _
  by_cases hf : ds_faithful phi
  · left; exact hf
  · right
    unfold ds_faithful at hf; push_neg at hf
    rcases hf with ⟨a, b, hab⟩
    exact ⟨a, b, by unfold translationLoss;
      exact abs_pos.mpr (sub_ne_zero.mpr hab)⟩
\end{lstlisting}


% ══════════════════════════════════════════════════════════════════════
% Chapter 3: Defamiliarization Fatigue
% ══════════════════════════════════════════════════════════════════════

\chapter{Defamiliarization Fatigue: Novelty Gets Old}
\label{ch:defamiliarization}

\epigraph{``Art exists that one may recover the sensation of life;
it exists to make one feel things, to make the stone stony.''}{Viktor Shklovsky, \textit{Art as Technique} (1917)}

Shklovsky's \emph{defamiliarization} (\textit{ostranenie})---the technique
of making the familiar strange---is central to modernist aesthetics.
But does defamiliarization have diminishing returns?  Can you keep
``making strange'' with undiminished effect?

In this chapter, we prove that the answer is \textbf{no}.  The cocycle
condition forces each successive defamiliarization to relate
algebraically to the previous one.  The artist cannot simply ``keep doing
the same weird thing'' and expect the same effect.  Novelty is
algebraically constrained.

\section{Iterated Defamiliarization}

\begin{definition}[Iterated Defamiliarization]
Applying the same defamiliarizing context $c$ to a sign $a$
repeatedly:
\[
  D^0(a, c) = a, \qquad D^{n+1}(a, c) = D^n(a, c) \compose c.
\]
\end{definition}

\begin{lstlisting}
def iteratedDefamiliarize (a ctx : I) : N → I
  | 0 => a
  | n + 1 => compose (iteratedDefamiliarize a ctx n) ctx
\end{lstlisting}

\begin{theorem}[Weight Grows Monotonically]
\label{thm:weight-mono}
Each iteration of defamiliarization strictly increases or preserves weight:
\[
  \rs(D^{n+1}(a,c), D^{n+1}(a,c)) \geq \rs(D^n(a,c), D^n(a,c)).
\]
\end{theorem}

\begin{proof}
$D^{n+1}(a,c) = D^n(a,c) \compose c$, so by axiom E4:
$\rs(D^n(a,c) \compose c, D^n(a,c) \compose c) \geq \rs(D^n(a,c), D^n(a,c))$.
\end{proof}

\begin{lstlisting}
theorem defamiliarize_weight_grows (a ctx : I) (n : N) :
    rs (iteratedDefamiliarize a ctx (n + 1))
       (iteratedDefamiliarize a ctx (n + 1)) ≥
    rs (iteratedDefamiliarize a ctx n)
       (iteratedDefamiliarize a ctx n) := by
  simp only [iteratedDefamiliarize]
  exact compose_enriches' (iteratedDefamiliarize a ctx n) ctx
\end{lstlisting}

\section{The Cocycle Constraint on Successive Novelty}

\begin{theorem}[Defamiliarization Cocycle]
\label{thm:defam-cocycle}
The emergence at the second defamiliarization is constrained by the
emergence at the first:
\[
  \emerge(a, c, \mathrm{obs}) + \emerge(a \compose c, c, \mathrm{obs})
  = \emerge(c, c, \mathrm{obs}) + \emerge(a, c \compose c, \mathrm{obs}).
\]
\end{theorem}

\begin{proof}
This is the cocycle condition with $b = c_1 = c$ (the same context
applied twice).
\end{proof}

\begin{lstlisting}
theorem defamiliarization_cocycle (a ctx obs : I) :
    emergence a ctx obs +
    emergence (compose a ctx) ctx obs =
    emergence ctx ctx obs +
    emergence a (compose ctx ctx) obs :=
  cocycle_condition a ctx ctx obs
\end{lstlisting}

\begin{remark}[Why this is surprising]
This is the key result.  It says that the emergence of the \emph{second}
application of context $c$ is \emph{not independent} of the first.
Rearranging:
\[
  \emerge(a \compose c, c, \mathrm{obs}) =
  \emerge(c, c, \mathrm{obs}) +
  \emerge(a, c \compose c, \mathrm{obs}) -
  \emerge(a, c, \mathrm{obs}).
\]
The second defamiliarization's effect is completely determined by three
quantities: (1) the self-emergence of $c$ (the context's own novelty),
(2) the emergence of $a$ with the doubled context $c \compose c$, and
(3) the emergence of the first application.  The artist cannot freely
choose the second novelty effect---it is constrained.

In particular, if $c$ has zero self-emergence ($\emerge(c,c,\mathrm{obs}) = 0$)
and $a$'s emergence with $c \compose c$ equals its emergence with $c$,
then the second application has \emph{exactly the same} effect as the
first.  And the third application is then constrained by the second.
The diminishing returns are structural, not accidental.
\end{remark}

\section{The Habituation Theorem}

\begin{theorem}[Habituation Constraint]
\label{thm:habituation}
After repeated defamiliarization with the same context, the emergence
at each step is fully determined by earlier steps:
\[
  \emerge(a \compose c, c, \mathrm{obs})
  = \emerge(c, c, \mathrm{obs}) + \emerge(a, c^2, \mathrm{obs})
    - \emerge(a, c, \mathrm{obs}).
\]
\end{theorem}

\begin{proof}
Rearranging the cocycle condition.
\end{proof}

\begin{lstlisting}
theorem habituation_constraint (a ctx obs : I) :
    emergence (compose a ctx) ctx obs =
    emergence ctx ctx obs +
    emergence a (compose ctx ctx) obs -
    emergence a ctx obs := by
  linarith [cocycle_condition a ctx ctx obs]
\end{lstlisting}

\begin{interpretation}[The algebra of artistic fatigue]
This theorem formalizes what every artist knows intuitively: you can't
keep doing the same trick.  The effect of the $n$-th weird thing is
algebraically determined by the $(n{-}1)$-th.  Each new iteration must
satisfy the cocycle constraint, which limits the freedom of the artist.

This is not merely a bound---it is an \emph{exact equation}.  The
emergence at step $n+1$ is a linear function of the emergence at step $n$,
plus corrections from the context's self-interaction and the
``accumulated'' emergence.  The artist's freedom is \emph{one degree of
freedom less} than naive independence would suggest.

In the language of cohomology, defamiliarization defines a 1-cochain,
and the cocycle condition says this cochain must be a cocycle.  The
artist is constrained to move within a cohomology class.
\end{interpretation}

\section{Void Defamiliarization Does Nothing}

\begin{theorem}[Void Defamiliarization Is Trivial]
\label{thm:void-defam}
Using silence as a defamiliarizing context has zero effect:
\[
  D^n(a, \void) = a \quad \text{for all } n.
\]
\end{theorem}

\begin{proof}
By induction on $n$.  Base case: $D^0(a, \void) = a$.
Inductive step: $D^{n+1}(a, \void) = D^n(a, \void) \compose \void
= a \compose \void = a$.
\end{proof}

\begin{lstlisting}
theorem void_defamiliarization (a obs : I) (n : N) :
    iteratedDefamiliarize a (void : I) n = a := by
  induction n with
  | zero => rfl
  | succ n ih => simp [iteratedDefamiliarize, ih, void_right']
\end{lstlisting}

\section{The Emergence Ratio Shrinks}

\begin{theorem}[Emergence-Weight Ratio Bound]
\label{thm:ratio-bound}
At each step, the emergence is bounded by
\[
  \emerge(a, c, \mathrm{obs})^2 \leq \rs(a \compose c, a \compose c) \cdot \rs(\mathrm{obs}, \mathrm{obs}).
\]
Since the left-hand side involves $\rs(a \compose c, a \compose c)$
which grows monotonically (Theorem~\ref{thm:weight-mono}), the
\emph{ratio} $\emerge^2 / \rs(D^n, D^n)$ can only decrease as the
weight grows.
\end{theorem}

\begin{proof}
Axiom E3 (emergence bounded).
\end{proof}

\begin{remark}[Relative novelty diminishes]
Even if the absolute emergence stays constant (a theoretical possibility),
the \emph{relative} novelty---emergence as a fraction of total weight---must
decrease.  This is the diminishing-returns theorem: each iteration of
defamiliarization contributes proportionally less to the total meaning
of the accumulated sign chain.  The $100$th repetition of ``make it
strange'' adds less \emph{relative} novelty than the first.
\end{remark}


% ══════════════════════════════════════════════════════════════════════
% Chapter 4: The Kitsch Theorem
% ══════════════════════════════════════════════════════════════════════

\chapter{The Kitsch Theorem: When Beauty Destroys Beauty}
\label{ch:kitsch}

\epigraph{``Kitsch causes two tears to flow in quick succession.
The first tear says: How nice to see children running on the grass!
The second tear says: How nice to be moved, together with all mankind,
by children running on the grass!''}{Milan Kundera, \textit{The Unbearable Lightness of Being}}

\section{Aesthetic Non-Monotonicity}

The most counter-intuitive result in our entire theory may be this:
\textbf{adding a beautiful element to a beautiful composition can make
the total aesthetic impact worse}.  In the language of economics, aesthetic
value is not a monotone function of its inputs.

\begin{definition}[Total Aesthetic Impact]
\[
  \mathrm{impact}(a, b, \mathrm{obs}) := \rs(a \compose b, \mathrm{obs}).
\]
\end{definition}

\begin{theorem}[The Kitsch Theorem]
\label{thm:kitsch}
\[
  \mathrm{impact}(a, b, \mathrm{obs})
  = \rs(a, \mathrm{obs}) + \rs(b, \mathrm{obs}) + \emerge(a, b, \mathrm{obs}).
\]
When $\emerge(a, b, \mathrm{obs}) < 0$, the total impact is \emph{less}
than the sum of individual impacts.
\end{theorem}

\begin{proof}
The resonance composition equation gives the decomposition directly.
\end{proof}

\begin{lstlisting}
theorem kitsch_theorem (a b observer : I) :
    totalAestheticImpact a b observer =
    rs a observer + rs b observer +
    emergence a b observer := by
  unfold totalAestheticImpact; rw [rs_compose_eq]
\end{lstlisting}

\begin{remark}[Why this is surprising]
This overturns the ``more is more'' intuition in aesthetics.  Consider:
\begin{itemize}
\item A beautifully written novel with an unnecessarily beautiful
  epilogue that undermines the ambiguity of the ending.
\item A perfectly composed photograph with an added Instagram filter
  that makes it look cheap.
\item A Baroque church with one too many gold cherubs.
\end{itemize}
In each case, the added element has positive resonance on its own, but
its \emph{emergence} with the existing work is negative.  The
composition clashes.  The whole is less than the sum of its parts.

This is the mathematical definition of \textbf{kitsch}: art where the
emergence is zero or negative, where addition is purely additive or
actively destructive.  The kitsch sensibility adds more beauty without
concern for how it interacts with what's already there.
\end{remark}

\section{The Over-Ornamentation Bound}

\begin{theorem}[Over-Ornamentation]
\label{thm:over-ornament}
If emergence is negative ($\emerge(a, b, \mathrm{obs}) < 0$), then
\[
  \mathrm{impact}(a, b, \mathrm{obs}) < \rs(a, \mathrm{obs}) + \rs(b, \mathrm{obs}).
\]
Adding $b$ makes the total impact \emph{strictly worse} than having
both parts separately.
\end{theorem}

\begin{proof}
Follows immediately from the Kitsch Theorem and the sign of the
emergence term.
\end{proof}

\begin{lstlisting}
theorem over_ornamentation (a b observer : I)
    (h : emergence a b observer < 0) :
    totalAestheticImpact a b observer <
    rs a observer + rs b observer := by
  unfold totalAestheticImpact
  linarith [rs_compose_eq a b observer]
\end{lstlisting}

\section{The Goldilocks Principle}

\begin{theorem}[The Goldilocks Principle]
\label{thm:goldilocks}
For any composition $a \compose b$:
\begin{enumerate}
\item If $\emerge(a, b, \mathrm{obs}) = 0$: impact equals the sum of parts (neutral).
\item If $\emerge(a, b, \mathrm{obs}) > 0$: impact exceeds the sum (synergy).
\item If $\emerge(a, b, \mathrm{obs}) < 0$: impact is less than the sum (clash).
\end{enumerate}
\end{theorem}

\begin{proof}
All three follow from the Kitsch Theorem by comparing signs.
\end{proof}

\begin{lstlisting}
theorem goldilocks_principle (a b obs : I) :
    (emergence a b obs = 0 →
      totalAestheticImpact a b obs = rs a obs + rs b obs) ∧
    (emergence a b obs > 0 →
      totalAestheticImpact a b obs > rs a obs + rs b obs) ∧
    (emergence a b obs < 0 →
      totalAestheticImpact a b obs < rs a obs + rs b obs) := by
  unfold totalAestheticImpact
  refine ⟨?_, ?_, ?_⟩ <;> intro h <;>
    linarith [rs_compose_eq a b obs]
\end{lstlisting}

\begin{interpretation}[The aesthetics of emergence]
The Goldilocks Principle says that aesthetic composition is not about
the quality of individual elements but about the \emph{emergence} between
them.  A mediocre element with positive emergence contributes more than
a brilliant element with negative emergence.  The art is in the
combination, not the parts.

This is the mathematical content of the old clich\'e ``less is more'':
removing an element with negative emergence \emph{increases} the total
impact.  The Goldilocks Principle gives the exact criterion for when
``more is more,'' ``more is less,'' and ``more is just more.''
\end{interpretation}

\section{Double Kitsch}

\begin{theorem}[Double Kitsch Is Predictable]
\label{thm:double-kitsch}
If $a$ is left-linear (kitsch), then
\[
  \mathrm{impact}(a, b, \mathrm{obs}) = \rs(a, \mathrm{obs}) + \rs(b, \mathrm{obs}).
\]
Kitsch produces exactly the sum of parts---no surprise, no synergy,
no clash.
\end{theorem}

\begin{proof}
Left-linearity means $\emerge(a, b, c) = 0$ for all $b, c$.
The Kitsch Theorem then gives the result.
\end{proof}

\begin{lstlisting}
theorem double_kitsch (a b observer : I)
    (ha : leftLinear a) :
    totalAestheticImpact a b observer =
    rs a observer + rs b observer := by
  unfold totalAestheticImpact
  rw [rs_compose_eq]; linarith [ha b observer]
\end{lstlisting}

\section{Weight Grows But Resonance Can Shrink}

\begin{theorem}[The Over-Production Paradox]
\label{thm:over-production}
The self-resonance (weight) of a composition always exceeds that of its
first element, but the resonance with a \emph{specific observer} can
decrease:
\[
  \rs(a \compose b, a \compose b) \geq \rs(a, a)
\]
but it is possible that $\rs(a \compose b, \mathrm{obs}) < \rs(a, \mathrm{obs})$.
\end{theorem}

\begin{proof}
The first inequality is axiom E4.  For the second, when
$\rs(b, \mathrm{obs}) + \emerge(a, b, \mathrm{obs}) < 0$,
we have $\rs(a \compose b, \mathrm{obs}) < \rs(a, \mathrm{obs})$.
\end{proof}

\begin{remark}[Why this is surprising]
This is the paradox of the over-produced film.  A film with a
\$200 million budget (high weight) can be less impactful to audiences
(less resonance with the observer) than a \$2 million indie film.
Weight measures \emph{quantity of stuff}; resonance measures
\emph{impact on a particular audience}.  They are different things,
and one can increase while the other decreases.
\end{remark}


% ══════════════════════════════════════════════════════════════════════
% Chapter 5: Narrative Order Dominates Content
% ══════════════════════════════════════════════════════════════════════

\chapter{Narrative Order Dominates Content}
\label{ch:order}

\epigraph{``The king died and then the queen died'' is a story.
``The king died, and then the queen died of grief'' is a plot.}{E.\ M.\ Forster, \textit{Aspects of the Novel}}

\section{The Order Sensitivity Theorem}

One of the most striking results in narrative theory: the \emph{order}
in which elements appear can change the meaning more than replacing
an element entirely.

\begin{definition}[Narrative Order Effect]
The \emph{order effect} of swapping $a$ and $b$ in a narrative:
\[
  \mathrm{order}(a, b, \mathrm{obs}) :=
  \rs(a \compose b, \mathrm{obs}) - \rs(b \compose a, \mathrm{obs}).
\]
\end{definition}

\begin{theorem}[Order Effect Equals Emergence Difference]
\label{thm:order-emergence}
\[
  \mathrm{order}(a, b, \mathrm{obs}) =
  \emerge(a, b, \mathrm{obs}) - \emerge(b, a, \mathrm{obs}).
\]
\end{theorem}

\begin{proof}
Expanding both resonances using the composition equation:
\begin{align*}
  \rs(a \compose b, \mathrm{obs}) &= \rs(a, \mathrm{obs}) + \rs(b, \mathrm{obs})
    + \emerge(a, b, \mathrm{obs}) \\
  \rs(b \compose a, \mathrm{obs}) &= \rs(b, \mathrm{obs}) + \rs(a, \mathrm{obs})
    + \emerge(b, a, \mathrm{obs}).
\end{align*}
Subtracting:
$\mathrm{order}(a, b, \mathrm{obs}) = \emerge(a, b, \mathrm{obs}) - \emerge(b, a, \mathrm{obs})$.
\end{proof}

\begin{lstlisting}
theorem narrativeOrderEffect_eq (a b observer : I) :
    narrativeOrderEffect a b observer =
    emergence a b observer - emergence b a observer := by
  unfold narrativeOrderEffect emergence; ring
\end{lstlisting}

\begin{remark}[Why this is surprising]
The order effect depends \emph{only} on emergence---it is \emph{purely}
a nonlinear phenomenon.  The linear parts (individual resonances)
cancel exactly.  This means that in any ideatic space where emergence is
nonzero, order \emph{must} matter.  The noncommutative structure of
composition is not an accident but a \emph{consequence} of having
nontrivial emergence.

Compare this with the effect of \emph{replacing} element $a$ with $a'$:
\[
  \rs(a' \compose b, \mathrm{obs}) - \rs(a \compose b, \mathrm{obs})
  = (\rs(a', \mathrm{obs}) - \rs(a, \mathrm{obs}))
    + (\emerge(a', b, \mathrm{obs}) - \emerge(a, b, \mathrm{obs})).
\]
The content change involves \emph{both} a linear term (resonance difference)
\emph{and} an emergence term.  The order change involves \emph{only}
emergence.  When emergence dominates---when ideas interact strongly---order
dominates content.

This is the formal version of ``\emph{how} you tell the story matters
more than \emph{what} the story is about.''
\end{remark}

\section{Weight Is Order-Blind; Meaning Is Not}

\begin{theorem}[Three-Act Weight Invariance]
\label{thm:weight-invariance}
In a three-element narrative $(a, b, c)$, the total weight
(self-resonance) is invariant under re-bracketing:
\[
  \rs((a \compose b) \compose c, (a \compose b) \compose c)
  = \rs(a \compose (b \compose c), a \compose (b \compose c)).
\]
\end{theorem}

\begin{proof}
By associativity: $(a \compose b) \compose c = a \compose (b \compose c)$.
\end{proof}

\begin{lstlisting}
theorem three_act_weight_invariant (a b c : I) :
    rs (compose (compose a b) c)
       (compose (compose a b) c) =
    rs (compose a (compose b c))
       (compose a (compose b c)) := by
  rw [compose_assoc']
\end{lstlisting}

\begin{interpretation}[Weight vs.\ meaning]
A story's ``substance'' (weight) doesn't depend on how you group
the acts.  But its ``meaning'' (resonance with an external observer)
\emph{does}---because the observer is not part of the associative
rearrangement.  The internal structure is bracket-invariant, but the
external perception is not.

This formalizes a deep intuition in narratology: the same events,
arranged differently, produce different stories with different meanings,
even though they have the same ``material.''
\end{interpretation}

\section{The Cocycle Constrains Narrative Structure}

\begin{theorem}[Narrative Cocycle]
\label{thm:narrative-cocycle}
The emergence of a three-act narrative decomposes according to
the cocycle condition:
\[
  \emerge(a, b, \mathrm{obs}) + \emerge(a \compose b, c, \mathrm{obs})
  = \emerge(b, c, \mathrm{obs}) + \emerge(a, b \compose c, \mathrm{obs}).
\]
\end{theorem}

\begin{proof}
The cocycle condition, applied to narrative elements.
\end{proof}

\begin{lstlisting}
theorem order_determines_emergence (a b c observer : I) :
    emergence a b observer +
    emergence (compose a b) c observer =
    emergence b c observer +
    emergence a (compose b c) observer :=
  cocycle_condition a b c observer
\end{lstlisting}

\begin{remark}[Why this is surprising]
Two arrangements of the same three elements produce the \emph{same}
total emergence sum (the cocycle guarantees this), but \emph{different}
individual emergence terms.  The \emph{distribution} of meaning across
the narrative changes with order, even though the total is conserved.

This is the formal content of Propp's observation that the same narrative
functions (Villainy, Departure, Donor, etc.) can be arranged in different
orders to produce different tales---but the tales are not completely free.
They must satisfy the cocycle constraint.  Narrative structure has an
algebra, and that algebra has laws.
\end{remark}

\section{Self-Composition Has Zero Order Effect}

\begin{theorem}[Self-Order Invariance]
\label{thm:self-order}
Composing an element with itself produces zero order effect:
$\mathrm{order}(a, a, \mathrm{obs}) = 0$ for all observers.
\end{theorem}

\begin{proof}
$a \compose a = a \compose a$, so the difference is zero.
\end{proof}

\begin{lstlisting}
theorem narrativeOrderEffect_self (a observer : I) :
    narrativeOrderEffect a a observer = 0 := by
  unfold narrativeOrderEffect; ring
\end{lstlisting}

\begin{interpretation}[Repetition is symmetric]
Repeating the same element (anaphora) is order-insensitive.  This makes
formal sense: saying ``fire, fire'' is the same whether you parse it as
``fire then fire'' or ``fire then fire.''  The order effect is a measure
of \emph{asymmetry} between two \emph{different} elements.  Repetition,
by definition, has no asymmetry.
\end{interpretation}
